\pdfbookmark[1]{Abstract}{Abstract}
\chapter*{Resumen}
Este trabajo presenta el diseño, la simulación, la fabricación y la caracterización experimental de un amplificador de potencia Clase AB en banda S, destinado a aplicaciones de enlace descendente en CubeSats de órbita terrestre baja, junto con el diseño de un acoplador híbrido en cuadratura concebido como primera etapa hacia una futura arquitectura Doherty. El amplificador emplea el transistor discreto GaN HEMT CGH40006S, seleccionado por su elevada densidad de potencia y su tolerancia a las condiciones térmicas y de radiación propias del entorno orbital. Un análisis de radioenlace determinó que una potencia de salida de 1~W resulta suficiente para sostener un enlace descendente con modulación QPSK a una tasa de 2,96~Mbps para ángulos de elevación superiores a 20°. El circuito se implementó sobre un sustrato FR4 de cuatro capas, con redes de adaptación de impedancia realizadas mediante componentes distribuidos. Simulaciones de load-pull bajo análisis de balance armónico de gran señal permitieron determinar la impedancia de carga óptima que maximiza la eficiencia de agregado de potencia (PAE). La caracterización de pequeña señal del prototipo fabricado arrojó una frecuencia central de 2,11~GHz, una ganancia de 13,26~dB y un ancho de banda de 70~MHz. Las mediciones de gran señal arrojaron un punto de compresión de 1~dB (P1dB) de 30,3~dBm, un OIP3 de 39,83~dBm y una eficiencia de drain de 32,94\,\% para una potencia de salida de 30,92~dBm. Por su parte, el acoplador híbrido en cuadratura, implementado mediante una topología de línea ramificada plegada para reducir sus dimensiones, alcanzó experimentalmente una frecuencia central de 2,019~GHz, con un desfasaje de 92,57° entre los puertos de salida en la frecuencia de interés. La concordancia obtenida entre los resultados simulados y medidos de ambos circuitos valida la metodología de diseño empleada, consolidando una plataforma de desarrollo y caracterización de módulos de radiofrecuencia reutilizable para futuros desarrollos aeroespaciales.

\keywordses{Amplificador de potencia, GaN, CubeSat, Comunicación satelital, Banda S, Acoplador híbrido en cuadratura.}

\MediaOptionLogicBlank